\begin{document}

\title{原始 HTML 与自定义组件}
\author{MetaBlog 示例}{}{example@example.com}

\begin{abstract}
这篇文章展示可信原始 HTML 块、外部 HTML 片段导入、行内原始 HTML，以及站点级自定义组件。
\end{abstract}

\section{原始 HTML 块}

当可信 HTML 片段需要直接进入页面时，可以使用 \texttt{html} 环境。

\begin{html}
<div style="border-left: 6px solid #2bb7b3; padding: 12px 16px; background: #f7fbfd;">
  <strong>原始 HTML 块：</strong>这段标记会不经过 escape 直接输出。
</div>
\end{html}

\section{导入 HTML 文件}

使用 \verb|\importHTML{...}| 可以加载相对于主 TeX 文件的可信 HTML 文件。

\importHTML{partials/profile-card.html}

兼容别名 \verb|\inputHTML{...}| 遵循同样规则。它和 \verb|\input{...}| 不同，专门用于导入原始 HTML 片段。

\section{行内原始 HTML}

如果前后没有空白行，原始 HTML 可以作为段落内部内容出现：
\begin{html}<span style="color:#256d85;font-weight:700;">这是一段行内 HTML 文本</span>\end{html}
并继续留在同一个段落中。

如果使用 \texttt{<div>} 这类块级 HTML，应当用空白行把它和上下文分开，避免浏览器自动闭合段落。

\section{自定义组件}

这个示例站点包含两个自定义组件：

\begin{itemize}
  \item \texttt{data/custom_components/page_footing.tex}，插入到每个页面底部。
  \item \texttt{data/custom_components/article_stat.tex}，插入到每篇文章的元信息下方。
\end{itemize}

组件文件只需要写内部内容。MetaBlog 会自动为它们包裹正确的外层元素。

\end{document}
